% future/formalregress.tex
% mainfile: ../perfbook.tex
% SPDX-License-Identifier: CC-BY-SA-3.0

\section{形式化回归测试？}
\label{sec:future:Formal Regression Testing?}
%
\epigraph{没有实验的理论：
	  我们走得太远了吗？}
	 {Michael Mitzenmacher}

形式化验证长期以来在许多生产环境中证明了其
有用性~\cite{JamesRLarus2004RightingSoftware,AlBessey2010BillionLoCLater,ByronCook2018FormalAmazon,CaitlinSadowski2018staticAnalysisGoogle,DinoDistefano2019FBstaticAnalysis}。
然而，硬核形式化验证是否能被纳入复杂并发代码库（如Linux内核）
的持续集成中所使用的自动化回归测试套件，仍是一个悬而未决的问题。
虽然已经有了Linux内核SRCU的概念验证~\cite{LanceRoy2017CBMC-SRCU}，
但该测试针对的是最简单的RCU实现之一的一小部分，
而且已经证明很难使其跟上不断变化的Linux内核。
因此，值得探讨将形式化验证作为Linux内核回归测试的一等成员
所需的条件。

以下列表是一个好的
起点~\cite[slide 34]{PaulEMcKenney2015DagstuhlVerification}：

\begin{enumerate}
\item	任何所需的转换必须是自动化的。
\item	环境（包括内存排序）必须被正确处理。
\item	内存和CPU开销必须在可接受的适度范围内。
\item	必须提供指向bug位置的具体信息。
\item	除源代码和输入之外的信息范围必须适度。
\item	所定位的bug必须与代码用户相关。
\end{enumerate}

这个列表基于但比Richard Bornat的格言更为温和：
``形式化验证研究人员应该验证开发人员编写的代码，
使用他们编写时所用的语言，在代码运行的环境中运行，
在他们编写时就进行验证。''
以下各节讨论上述每项要求，
然后有一节展示几个工具在这些要求方面表现的记分卡。

\QuickQuiz{
	这个列表太乌托邦了！
	为什么不坚持形式化验证的现有技术水平？
}\QuickQuizAnswer{
	你有权对什么是乌托邦、什么不是乌托邦发表意见，
	但比起任何可能反对这些条目的人，
	我会更关注那些真正在这些条目上取得进展的人。
	这可能与我长期以来的经验有关，
	即有人试图说服我放弃其所偏爱的工具无法处理的特定事项。

	与此同时，请随意阅读那些真正取得进展的人所写的论文，
	例如这篇~\cite{DinoDistefano2019FBstaticAnalysis}。
}\QuickQuizEnd

\subsection{自动转换}
\label{sec:future:Automatic Translation}

虽然Promela和\co{spin}是无价的设计辅助工具，
但如果你需要对C语言程序进行形式化回归测试，
你必须在每次想要重新验证代码时手动将其翻译为Promela。
如果你的代码恰好在Linux内核中（每60--90天发布一次），
你将需要每年手动翻译四到六次。
随着时间推移，人为错误将悄然出现，
这意味着验证将不再与源代码匹配，使验证变得毫无用处。
重复验证显然要求形式化验证工具直接输入你的代码，
或者存在无bug的自动转换将你的代码转换为验证所需的形式。

PPCMEM和\co{herd}理论上可以直接输入汇编语言和C++代码，
但这些工具只能处理非常小的litmus测试，
这通常意味着你必须手动提取你的机制的核心部分。
与Promela和\co{spin}一样，PPCMEM和\co{herd}都极为有用，
但它们不太适合回归测试套件。

相比之下，\IXaltacr{\co{cbmc}}{cbmc}和\IX{Nidhugg}可以输入
合理大小（虽然仍然相当有限）的C程序，
如果它们的能力继续增长，完全可以成为回归测试套件的出色补充。
Coverity静态分析工具也输入C程序，而且可以处理非常大的程序，
包括Linux内核。
当然，Coverity的静态分析与\co{cbmc}和Nidhugg相比要简单得多。
另一方面，Coverity拥有一个包罗万象的``C程序''定义，
这带来了特殊的挑战~\cite{AlBessey2010BillionLoCLater}。
Amazon Web Services使用各种形式化验证工具，
包括\co{cbmc}，并将其中一些工具应用于
回归测试~\cite{ByronCook2018FormalAmazon}。
Google在大型Java代码库上直接使用许多相对简单的静态分析工具，
这些代码库可以说比C代码库更为同质~\cite{CaitlinSadowski2018staticAnalysisGoogle}。
Facebook对其代码库使用更激进的形式化验证形式，
包括并发分析~\cite{DinoDistefano2019FBstaticAnalysis,PeterWOHearn2019incorrectnessLogic}，
尽管尚未应用于Linux内核。
最后，Microsoft长期以来一直对其代码库使用
静态分析~\cite{JamesRLarus2004RightingSoftware}。

鉴于这份清单，显然可以创建直接消费生产级别源代码的
复杂形式化验证工具。

然而，以C代码作为输入的一个不足之处在于它假设编译器是正确的。
另一种方法是将C编译器生成的二进制文件作为输入，
从而考虑到任何相关的编译器bug。
这种方法已在许多验证工作中使用，
最著名的可能是seL4项目~\cite{ThomasSewell2013L4binaryVerification}。

\QuickQuiz{
	鉴于seL4项目中使用的各种验证器的开创性，
	为什么本章不更深入地介绍它们？
}\QuickQuizAnswer{
	毫无疑问，seL4项目使用的验证器能力相当强。
	然而，seL4最初是一个单CPU项目。
	虽然seL4已经获得了多处理器能力，
	但它目前使用的是非常粗粒度的锁，
	类似于Linux内核的旧Big Kernel Lock (BKL)\@。
	希望有一天将seL4的验证器添加到一本关于并行编程的书中
	是有意义的，但现在还不是时候。
}\QuickQuizEnd

然而，直接从源代码或二进制文件进行验证都具有
消除人工翻译错误的优势，
这对于可靠的回归测试至关重要。

这并不是说具有专用语言的工具是无用的。
相反，它们对于设计时验证非常有帮助，
正如在\cref{chp:Formal Verification}中讨论的那样。
然而，这类工具对于自动化回归测试并不特别有帮助，
而这事实上正是本节的主题。

\subsection{环境}
\label{sec:future:Environment}

形式化验证工具正确建模其环境至关重要。
一个常见的遗漏是\IX{memory model}，
其中大量形式化验证工具（包括Promela/\co{spin}）
被限制为\IXalth{sequential consistency}{sequential}{memory consistency}。
\cref{sec:formal:Is QRCU Really Correct?}中叙述的QRCU经历
是一个重要的警示故事。

Promela和\co{spin}假设顺序一致性，
这与现代计算机系统不太匹配，
正如在\cref{chp:Advanced Synchronization: Memory Ordering}中所见。
相比之下，PPCMEM和\co{herd}的一大优势是
它们对各种CPU家族内存模型的详细建模，
包括x86、\ARM、Power，以及在\co{herd}的情况下，
还有Linux内核内存模型~\cite{Alglave:2018:FSC:3173162.3177156}，
该模型已被接受进入Linux内核v4.17版本。

\co{cbmc}和Nidhugg工具提供了一些选择内存模型的能力，
但没有提供PPCMEM和\co{herd}那样的多样性。
然而，随着时间推移，更大规模的工具可能会采用更多种类的内存模型。

从更长远来看，如果形式化验证工具能够包含
I/O~\cite{PaulEMcKenney2016LinuxKernelMMIO}会很有帮助，
但这可能还需要一段时间才能实现。

尽管如此，无法匹配环境的工具仍然可以是有用的。
例如，许多并发bug在神话般的顺序一致系统上仍然是bug，
这些bug可以被用顺序一致性过度近似系统内存模型的
工具所发现。
然而，这些工具将无法找到涉及缺少内存排序指令的bug，
正如前述\cref{sec:formal:Is QRCU Really Correct?}
的警示故事中所指出的那样。

\subsection{开销}
\label{sec:future:Overhead}

几乎所有硬核形式化验证工具本质上都是指数级的，
这可能看起来令人沮丧，直到你考虑到许多最有趣的
软件问题实际上是不可判定的。
然而，即使在指数级之间也存在程度差异。

PPCMEM在设计上是未优化的，
以便更好地保证所关注的内存模型被准确表示。
\co{herd}工具则更激进地优化，
如\cref{sec:formal:Axiomatic Approaches}中所述，
因此比PPCMEM快了数个数量级。
然而，PPCMEM和\co{herd}都针对的是非常小的litmus测试，
而不是较大的代码体。

相比之下，Promela/\co{spin}、\co{cbmc}和Nidhugg是为
（某种程度上）更大的代码体设计的。
Promela/\co{spin}被用于验证好奇号火星车的
文件系统~\cite{DBLP:journals/amai/GroceHHJX14}，
如前所述，\co{cbmc}和Nidhugg都被应用于Linux内核RCU\@。

如果启发式方法的进步继续保持过去三十年的速度，
我们可以期待形式化验证的开销大幅减少。
也就是说，组合爆炸仍然是组合爆炸，
这将预期会严格限制可以被验证的程序的大小，
无论启发式方法是否继续改进。

然而，组合爆炸的另一面是马其顿的腓力二世的
永恒建议：``分而治之。''
如果一个大程序可以被分割并验证各个部分，
结果可以是组合\emph{内爆}~\cite{PaulEMcKenney2011Verico}。
一个自然的分割点是在API边界上，例如锁原语的边界。
一次验证通过可以验证锁实现是正确的，
额外的验证通过可以验证锁API的正确使用。

\begin{listing}
\input{CodeSamples/formal/herd/C-SB+l-o-o-u+l-o-o-u-C@whole.fcv}
\caption{用\tco{cmpxchg_acquire()}模拟锁}
\label{lst:future:Emulating Locking with cmpxchg}
\end{listing}

\begin{table}
\rowcolors{1}{}{lightgray}
\renewcommand*{\arraystretch}{1.1}
\small
\centering
\begin{tabular}{S[table-format=1.0]S[table-format=1.3]S[table-format=2.3]}
	\toprule
	\multicolumn{1}{c}{\# 线程数} & \multicolumn{1}{c}{锁} &
			\multicolumn{1}{c}{\tco{cmpxchg_acquire}} \\
	\midrule
	2 & 0.004 &  0.022 \\
	3 & 0.041 &  0.743 \\
	4 & 0.374 & 59.565 \\
	5 & 4.905 &        \\
	\bottomrule
\end{tabular}
\caption{模拟锁：
			    性能（秒）}
\label{tab:future:Emulating Locking: Performance (s)}
\end{table}

这种方法的性能优势可以使用Linux内核
内存模型~\cite{Alglave:2018:FSC:3173162.3177156}来演示。
该模型提供了\co{spin_lock()}和\co{spin_unlock()}原语，
但这些原语也可以使用\co{cmpxchg_acquire()}和
\co{smp_store_release()}来模拟，
如\cref{lst:future:Emulating Locking with cmpxchg}所示
（\path{C-SB+l-o-o-u+l-o-o-*u.litmus}和\path{C-SB+l-o-o-u+l-o-o-u*-C.litmus}）。
\Cref{tab:future:Emulating Locking: Performance (s)}
比较了使用模型的\co{spin_lock()}和\co{spin_unlock()}
与列表中所示的模拟这些原语的性能和可扩展性。
差异是显著的：
在四个进程时，模型比模拟快两个数量级以上！

\QuickQuiz{
\begin{fcvref}[ln:future:formalregress:C-SB+l-o-o-u+l-o-o-u-C:whole]
	为什么要在\cref{lst:future:Emulating Locking with cmpxchg}的
	\clnref{filter_}上使用单独的\co{filter}命令，
	而不是直接将条件添加到\co{exists}子句中？
	而且使用\co{xchg_acquire()}代替\co{cmpxchg_acquire()}
	不是更简单吗？
\end{fcvref}
}\QuickQuizAnswer{
	\co{filter}子句使得\co{herd}工具在比\co{exists}子句
	更早的处理阶段丢弃执行，这提供了显著的加速。

\begin{table}
\rowcolors{7}{lightgray}{}
\renewcommand*{\arraystretch}{1.1}
\small
\centering
\begin{tabular}{S[table-format=1.0]S[table-format=1.3]S[table-format=2.3]
		S[table-format=3.3]S[table-format=2.3]S[table-format=3.3]}
	\toprule
	& & \multicolumn{2}{c}{\tco{cmpxchg_acquire()}}
		& \multicolumn{2}{c}{\tco{xchg_acquire()}} \\
	\cmidrule(l){3-4} \cmidrule(l){5-6}
	\multicolumn{1}{c}{\#} & \multicolumn{1}{c}{锁}
		& \multicolumn{1}{c}{\tco{filter}}
			& \multicolumn{1}{c}{\tco{exists}}
				& \multicolumn{1}{c}{\tco{filter}}
					& \multicolumn{1}{c}{\tco{exists}} \\
	\cmidrule{1-1} \cmidrule(l){2-2} \cmidrule(l){3-4} \cmidrule(l){5-6}
	2 & 0.004 &  0.022 &   0.039 &  0.027 &  0.058 \\
	3 & 0.041 &  0.743 &   1.653 &  0.968 &  3.203 \\
	4 & 0.374 & 59.565 & 151.962 & 74.818 & 500.96 \\
	5 & 4.905 &        &         &        &        \\
	\bottomrule
\end{tabular}
\caption{模拟锁：
			    性能比较（秒）}
\label{tab:future:Emulating Locking: Performance Comparison (s)}
\end{table}

	至于\co{xchg_acquire()}，这个原子操作无论锁获取是否成功
	都会执行写操作，这意味着使用\co{xchg_acquire()}的模型
	将比使用\co{cmpxchg_acquire()}的模型有更多操作，
	后者在获取失败的情况下不会执行写操作。
	更多的写操作意味着更多的组合爆炸，
	如\cref{tab:future:Emulating Locking: Performance Comparison (s)}所示
	（\path{C-SB+l-o-o-u+l-o-o-*u.litmus}、
	\path{C-SB+l-o-o-u+l-o-o-u*-C.litmus}、
	\path{C-SB+l-o-o-u+l-o-o-u*-CE.litmus}、
	\path{C-SB+l-o-o-u+l-o-o-u*-X.litmus}和
	\path{C-SB+l-o-o-u+l-o-o-u*-XE.litmus}）。
	该表清楚地表明\co{cmpxchg_acquire()}
	优于\co{xchg_acquire()}，
	而使用\co{filter}子句优于使用\co{exists}子句。
}\QuickQuizEnd

如果工具能够自动将大程序分割、验证各个部分，
然后验证部分的组合，那当然是非常有用的。
在此之前，大程序的验证将需要大量的人工干预。
这种干预最好通过脚本来中介，
以便更可靠地在每个版本发布时进行重复验证，
最好最终以适合持续集成的方式进行。
而Facebook的Infer工具已经朝着这个方向迈出了重要步伐，
通过组合性和抽象化~\cite{SamBlackshear2018RacerD,DinoDistefano2019FBstaticAnalysis}。

无论如何，我们可以期待形式化验证能力随时间继续增强，
任何这样的增强都将反过来提高形式化验证对回归测试的适用性。

\subsection{定位Bug}
\label{sec:future:Locate Bugs}

任何有一定规模的软件工件都包含bug。
因此，一个只报告bug存在与否的形式化验证工具并不特别有用。
需要的是一个至少提供\emph{一些}关于bug位置和性质信息的工具。

\co{cbmc}的输出包含映射回源代码的回溯信息，
类似于Promela/\co{spin}的回溯，Nidhugg也是如此。
当然，这些回溯可能相当长，分析它们可能相当乏味。
然而，这样做通常比用传统方式定位bug要快得多、愉快得多。

此外，形式化验证工具的最简单测试之一是bug注入。
毕竟，不仅我们中的任何人都可以写
\co{printf("VERIFIED\\n")}，而且事实是
形式化验证工具的开发人员和我们其他人一样容易犯错。
因此，仅仅宣称存在bug的形式化验证工具从根本上来说
不太值得信赖，因为更难在实际代码上对其进行验证。

撇开这些不谈，编写形式化验证工具的人可以利用现有工具。
例如，一个旨在仅确定严重但罕见的bug是否存在的工具
可以利用二分法。
如果被测程序的旧版本不包含该bug，
但新版本包含，那么可以使用二分法快速定位引入bug的提交，
这可能是找到并修复bug的足够信息。
当然，这种策略对于常见bug不太有效，
因为在这种情况下二分法会失败，
因为所有提交都至少有一个常见bug的实例。

因此，许多形式化验证工具提供的执行跟踪将继续是有价值的，
特别是对于复杂且难以理解的bug。
此外，最近的工作应用了\emph{不正确性逻辑}形式化方法，
类似于用于完整正确性证明的传统Hoare逻辑，
但其唯一目的是查找bug~\cite{PeterWOHearn2019incorrectnessLogic}。

\subsection{最小脚手架}
\label{sec:future:Minimal Scaffolding}

在过去，形式化验证研究人员要求提供完整的规范，
软件将根据该规范进行验证。
不幸的是，一个数学上严格的规范可能比实际代码还要大，
规范的每一行与代码的每一行一样可能包含bug。
一项证明代码忠实实现了规范的形式化验证工作
将是两者之间bug对bug兼容性的证明，
这可能不太有帮助。

更糟糕的是，许多软件工件的需求本质上是经验性的，
包括Linux内核的RCU~\cite{PaulEMcKenney2015RCUreqts1,PaulEMcKenney2015RCUreqts2,PaulEMcKenney2015RCUreqts3}。\footnote{
	或者，用形式化验证的术语来说，Linux内核的RCU具有
	\emph{不完整的规范}。}
对于这种常见类型的软件，完整的规范是一种礼貌的虚构。
完整的规范对于硬件同样是虚构的，
这一点在2017年末的Meltdown和Spectre侧信道攻击
中得到了充分证明~\cite{JannHorn2018MeltdownSpectre}。

这种情况可能会让人对实际软件和硬件工件的形式化验证
完全丧失希望，但事实证明可以做很多事情。
例如，设计和编码规则可以作为部分规范，
代码中包含的断言也是如此。
事实上，\co{cbmc}和Nidhugg等形式化验证工具都会检查
可以触发的断言，隐式地将这些断言作为规范的一部分。
然而，断言也是代码的一部分，这使得它们不太可能变得过时，
特别是如果代码也经过了压力测试。\footnote{
	你\emph{确实}对你的代码进行压力测试，对吧？}
\co{cbmc}工具还检查数组越界引用，
从而隐式地将其添加到规范中。
前述的不正确性逻辑也可以被认为使用了一个隐式的
无bug存在的规范~\cite{PeterWOHearn2019incorrectnessLogic}。

这种隐式规范方法很有道理，特别是如果你将形式化验证
不是视为完整的正确性证明，
而是视为一种具有与常见情况（即测试）不同优势和劣势的
替代验证形式。
从这个角度来看，软件总会有bug，因此任何有助于
找到这些bug的任何种类的工具都是非常好的。

\subsection{相关的Bug}
\label{sec:future:Relevant Bugs}

找到bug并修复它们当然是任何类型验证工作的全部意义。
显然，应避免误报。
但即使没有误报，bug和bug也是不同的。

例如，假设一个软件工件恰好有100个剩余bug，
每个bug平均每一百万年的运行时间才会出现一次。
进一步假设一个全知的形式化验证工具定位了所有100个bug，
开发人员勤勉地修复了它们。
这个软件工件的可靠性会发生什么变化？

答案是可靠性\emph{降低}了。

要理解这一点，请记住历史经验表明大约7\pct{}的修复
会引入新的bug~\cite{RexBlack2012SQA}。
因此，修复这100个bug（其组合平均故障间隔时间（MTBF）约为10,000年）
将引入另外七个bug。
历史统计表明每个新bug的MTBF远小于70,000年。
这反过来表明这七个新bug的组合MTBF
很可能远小于10,000年，
这又意味着善意修复原来100个bug的行为实际上降低了
整体软件的可靠性。

\QuickQuizSeries{%
\QuickQuizB{
	我们怎么知道已知bug的MTBF是对尚未发现的bug的
	MTBF的良好估计？
}\QuickQuizAnswerB{
	我们不知道，但这并不重要。

	要理解这一点，请注意7\pct{}的数字仅适用于
	随后被发现的注入bug：
	它必然忽略了任何从未被发现的注入bug。
	因此，已知bug的MTBF统计可能是随后被发现的注入bug的
	良好近似。

	本节的一个关键点是，我们应该更关心那些给用户带来不便的bug，
	而不是那些从未实际出现的其他bug。
	当然，这\emph{并非}说我们应该完全忽略尚未给用户带来
	不便的bug，只是我们应该适当地确定优先级，
	以便首先修复最重要和最紧急的bug。
}\QuickQuizEndB
%
\QuickQuizE{
	但形式化验证工具应该立即找到修复引入的所有bug，
	那么这为什么是个问题？
}\QuickQuizAnswerE{
	这是个问题，因为现实世界的形式化验证工具
	（与那些仅存在于形式化验证更激进倡导者的想象中的工具不同）
	不是全知的，因此只能定位某些类型的bug。
	仅举一个例子，形式化验证工具不太可能发现
	与遗漏断言对应的bug，或者等价地，
	与规范中未发现部分对应的bug。
}\QuickQuizEndE
}

更糟糕的是，想象另一个软件工件有一个bug平均每天失败一次，
还有99个每百万年失败一次的bug。
假设一个形式化验证工具定位了99个百万年的bug，
但未能找到那个一天一次的bug。
修复定位到的99个bug将花费时间和精力，降低可靠性，
而对那个可能正在造成尴尬甚至更糟糕后果的每天失败
毫无帮助。

因此，最好有一个能优先定位最麻烦bug的验证工具。
然而，如\cref{sec:future:Locate Bugs}中所述，
可以利用额外的工具。
一个强大的工具就是普通的测试。
了解了bug之后，应该可以为其构建特定的测试，
也可能使用\cref{sec:debugging:Hunting Heisenbugs}中描述的
一些技术来增加bug出现的概率。
这些技术应该允许计算bug原始故障率的粗略估计，
这反过来可以用于确定bug修复工作的优先级。

\QuickQuiz{
	但许多形式化验证工具一次只能找到一个bug，
	因此必须在工具能定位下一个bug之前修复每个bug。
	在这样的工具下如何确定bug修复工作的优先级？
}\QuickQuizAnswer{
	一种方法是提供一个可能不适合生产环境的简单修复，
	但允许工具定位下一个bug。
	另一种方法是限制配置或输入，
	使得到目前为止定位的bug不会发生。
	有许多类似的方法，但共同的主题是
	从工具的角度修复bug通常比构建和验证生产级别修复
	容易得多，关键是要优先考虑构建和验证
	生产级别修复所需的更大工作量。
}\QuickQuizEnd

最近有一些形式化验证工作优先考虑具有较少抢占的执行，
基于较少的抢占次数更可能发生的合理假设。

识别相关bug可能听起来要求太高了，
但如果我们要真正提高软件可靠性，这正是所需要的。

\subsection{形式化回归记分卡}
\label{sec:future:Formal Regression Scorecard}

\begin{table*}
% \rowcolors{6}{}{lightgray}
%\renewcommand*{\arraystretch}{1.1}
\small
\centering
\setlength{\tabcolsep}{2pt}
\begin{tabular}{lcccccccccc}
	\toprule
	& & Promela & & PPCMEM & & \tco{herd} & & \tco{cbmc} & & Nidhugg \\
	\midrule
	(1) 自动化 &
		& \cellcolor{red!50} &
			& \cellcolor{orange!50} &
				& \cellcolor{orange!50} &
					& \cellcolor{blue!50} &
						& \cellcolor{blue!50} \\
	\addlinespace[3pt]
	(2) 环境 &
		& \cellcolor{red!50} (MM) &
			& \cellcolor{green!50} &
				& \cellcolor{blue!50} &
					& \cellcolor{yellow!50} (MM) &
						& \cellcolor{orange!50} (MM) \\
	\addlinespace[3pt]
	(3) 开销 &
		& \cellcolor{yellow!50} &
			& \cellcolor{red!50} &
				& \cellcolor{yellow!50} &
					& \cellcolor{yellow!50} (SAT) &
						& \cellcolor{green!50} \\
	\addlinespace[3pt]
	(4) 定位Bug &
		& \cellcolor{yellow!50} &
			& \cellcolor{yellow!50} &
				& \cellcolor{yellow!50} &
					& \cellcolor{green!50} &
						& \cellcolor{green!50} \\
	\addlinespace[3pt]
	(5) 最小脚手架 &
		& \cellcolor{green!50} &
			& \cellcolor{yellow!50} &
				& \cellcolor{yellow!50} &
					& \cellcolor{blue!50} &
						& \cellcolor{blue!50} \\
	\addlinespace[3pt]
	(6) 相关Bug &
		& \cellcolor{yellow!50} ??? &
			& \cellcolor{yellow!50} ??? &
				& \cellcolor{yellow!50} ??? &
					& \cellcolor{yellow!50} ??? &
						& \cellcolor{yellow!50} ??? \\
	\bottomrule
\end{tabular}
\caption{形式化回归记分卡}
\label{tab:future:Formal Regression Scorecard}
\end{table*}

\Cref{tab:future:Formal Regression Scorecard}
显示了本章涵盖的形式化验证工具的粗略记分卡。
波长越短越好。

Promela需要手动翻译且仅支持
\IXalth{sequential consistency}{sequential}{memory consistency}，
因此其前两个单元格为红色。
它有合理的开销（就形式化验证而言）
且提供回溯信息，因此其接下来两个单元格为黄色。
尽管需要手动翻译，Promela以自然的方式处理断言，
因此其第五个单元格为绿色。

PPCMEM通常需要手动翻译，因为它支持的litmus测试规模很小，
因此其第一个单元格为橙色。
它处理多种\IX{memory model}，因此其第二个单元格为绿色。
它的开销相当高，因此其第三个单元格为红色。
它提供操作之间关系的图形显示，
虽不如回溯信息有帮助，但仍然相当有用，
因此其第四个单元格为黄色。
它需要构造\co{exists}子句且不能进行进程内断言，
因此其第五个单元格也是黄色。

\co{herd}工具具有与PPCMEM类似的大小限制，
因此\co{herd}的第一个单元格也是橙色。
它支持多种内存模型，因此其第二个单元格为蓝色。
它有合理的开销，因此其第三个单元格为黄色。
它的bug定位和断言能力与PPCMEM非常相似，
因此\co{herd}接下来两个单元格也是黄色。

\co{cbmc}工具直接输入C代码，因此其第一个单元格为蓝色。
它支持几种内存模型，因此其第二个单元格为黄色。
它有合理的开销，因此其第三个单元格也是黄色，
不过也许SAT求解器的性能将继续改善。
它提供回溯信息，因此其第四个单元格为绿色。
它直接从C代码中获取断言，因此其第五个单元格为蓝色。

Nidhugg也直接输入C代码，因此其第一个单元格也是蓝色。
它只支持几种内存模型，因此其第二个单元格为橙色。
它的开销相当低（就形式化验证而言），
因此其第三个单元格为绿色。
它提供回溯信息，因此其第四个单元格为绿色。
它直接从C代码中获取断言，因此其第五个单元格为蓝色。

那么第六行也就是最后一行呢？
现在判断这些工具在找到正确bug方面表现如何还为时过早，
因此它们都是带问号的黄色。

\QuickQuizSeries{%
\QuickQuizB{
	在\cref{tab:future:Formal Regression Scorecard}
	所示的记分卡中，测试会如何表现？
}\QuickQuizAnswerB{
	它将一路蓝色到底，
	可能的例外是第三行（开销），
	可能会因为测试难以找到不太可能出现的bug而被扣分。

	另一方面，不太可能出现的bug通常也是不相关的bug，
	因此你的结果可能会有所不同。

	很大程度上取决于你的安装基础的大小。
	如果你的代码只会在（比如说）10,000个系统上运行，
	Murphy实际上可以是一个非常好的人。
	所有可能出错的事情都会出错。
	最终。
	也许在地质时间尺度上。

	但如果你的代码运行在200亿个系统上，
	就像据说2017年底Linux内核那样，
	Murphy可以是一个真正的混蛋！
	所有可能出错的事情都会出错，
	而且可能出错得非常快！！！
}\QuickQuizEndB
%
\QuickQuizE{
	难道不是有比\cref{tab:future:Formal Regression Scorecard}
	中显示的更多的形式化验证系统吗？
}\QuickQuizAnswerE{
	确实如此！
	该表关注的是Paul使用过的那些，但其他一些也在证明其有用性。
	形式化验证在seL4项目中得到了大量
	使用~\cite{ThomasSewell2013L4binaryVerification}，
	其工具现在可以处理适度水平的并发。
	最近，Catalin Marinas使用Lamport的
	TLA工具~\cite{Lamport:2002:SST:579617}定位了Linux内核
	排队自旋锁（spinlock）实现中的一些前向进展bug。
	Will Deacon修复了这些bug~\cite{WillDeacon2018qspinlock}，
	而Catalin验证了Will的
	修复~\cite{CatalinMarinas2018qspinlockTLA}。

	较轻量级的形式化验证工具已在生产中得到大量
	使用~\cite{JamesRLarus2004RightingSoftware,AlBessey2010BillionLoCLater,ByronCook2018FormalAmazon,CaitlinSadowski2018staticAnalysisGoogle,DinoDistefano2019FBstaticAnalysis}。
}\QuickQuizEndE
}

再次请注意，该表对这些工具在回归测试中的使用进行评分。
仅仅因为它们中的许多不适合回归测试，
并不意味着它们是无用的，事实上，
它们中的许多已经多次证明了自己的价值。\footnote{
	仅举一个例子，Promela被用于验证好奇号火星车的文件系统。
	\emph{你的}形式化验证工具是否被用于目前在火星上运行的软件？？？}
只是不适合回归测试。

然而，这可能会改变。
毕竟，形式化验证工具在2010年代取得了令人印象深刻的进步。
如果这种进步继续下去，形式化验证可能会成为
并行程序员验证工具箱中不可或缺的工具。
